\documentclass[a4paper]{article}
\usepackage[affil-it]{authblk}
\usepackage[english]{babel}
\usepackage[backend=bibtex,style=numeric]{biblatex}

\usepackage{geometry}
\geometry{margin=1.5cm, vmargin={0pt,1cm}}
\setlength{\topmargin}{-1cm}
\setlength{\paperheight}{29.7cm}
\setlength{\textheight}{25.3cm}

% useful packages.
\usepackage{amsfonts}
\usepackage{amsmath}
\usepackage{amssymb}
\usepackage{amsthm}     
\usepackage{enumerate} %enumerate environment
\usepackage{graphicx}
\usepackage{multicol}
\usepackage{fancyhdr}
\usepackage{layout}
\usepackage{ctex}      %中文环境
\usepackage{booktabs}  %三线表
\usepackage{verbatim}  %comment environment
\usepackage{float}     %强制表格图片位置
\usepackage{listings}  %insert code in latex
\usepackage{fontspec}  %font support
\usepackage{tikz}
\usepackage{makecell}
\usetikzlibrary{intersections}
\usepackage{arydshln}

\usepackage[colorlinks,linkcolor=blue]{hyperref}
\usepackage{xcolor}
\definecolor{mygreen}{rgb}{0,0.6,0}
\definecolor{lightgrey}{rgb}{0.95,0.95,0.95}
\definecolor{winered}{rgb}{0.5,0,0}
\definecolor{mymauve}{rgb}{0.58,0,0.82}
\lstset{
      backgroundcolor=\color{lightgrey},
      basicstyle             = \ttfamily,
      breakatwhitespace      = false,
      breaklines             = true,
      captionpos             = none,
      commentstyle           = \color{mygreen}\bfseries,
      extendedchars          = false,
      keepspaces             = true,
      keywordstyle           = \color{winered},         % keyword style
      language               = C++,                     % the language of code
      otherkeywords          = {string},
      numberstyle            = \tiny\color{mygray},
      rulecolor              = \color{black},
      showspaces             = false,
      showstringspaces       = false,
      showtabs               = false,
      stepnumber             = 1,
      stringstyle            = \color{mymauve},         % string literal style
      tabsize                = 2,
      title                  = \lstname,
      aboveskip              = 0pt,
      belowskip              = 0pt,
      columns                = flexible
}

% some common command
\newcommand{\dif}{\mathrm{d}}
\newcommand{\avg}[1]{\left\langle #1 \right\rangle}
\newcommand{\norm}[1]{\left\| #1 \right\|}
\newcommand{\difFrac}[2]{\frac{\dif #1}{\dif #2}}
\newcommand{\pdfFrac}[2]{\frac{\partial #1}{\partial #2}}
\newcommand{\OFL}{\mathrm{OFL}}
\newcommand{\UFL}{\mathrm{UFL}}
\newcommand{\fl}{\mathrm{fl}}
\newcommand{\op}{\odot}
\newcommand{\Eabs}{E_{\mathrm{abs}}}
\newcommand{\Erel}{E_{\mathrm{rel}}}
\addbibresource{citation.bib}

\begin{document}
% =================================================
\title{Numerical Analysis chapter 6 theoretical homework}

\author{Zhouyue Qian 3220104074
  \thanks{Electronic address: \texttt{3220104074@edu.zju.cn}}}
\affil{(Information and Computational Science Class 2201), Zhejiang University }


\date{Due time: \today}

\maketitle

%\begin{abstract}
%    The abstract is not necessary for the theoretical homework, 
%    but for the programming project, 
%    you are encouraged to write one.      
%\end{abstract}

% ============================================
\section*{I}
\subsection*{a}
By Simpson's rule, we have $\int_{a}^{b} f(x) \,dx = \frac{b-a}{6}[f(a)+4f(\frac{b+a}{2})+f(b)]+R_s$.
\[
\begin{array}{c|cccc}
x & f(x) & f[x, x+1] & f[x, x+1, x+2] & f[x, x+1, x+2, x+3] \\
\hline
-1 & y(-1) & & & \\
0 & y(0) & y(0) -y(-1)  & & \\
0 & y(0) & y'(0) & y'(0) -y(0) +y(-1)& \\
1 & y(1) & y(1)-y(0) & y(1)-y(0) -y'(0)& \frac{1}{2}(y(1)-2y'(0)-y(-1))\\
\end{array}
\]
So,
\begin{align*}
  p_3(x) &= L_0+L_!(x+1)+L_2x(x+1)+L_3x^2(x+1)\\
  &=L_3x^3+(L_2+L_3)x^2+(L_1+L_2)x+L_0+L_1\\
  \int_{-1}^{1} p_3(x) \,dx & = \frac{2}{3}(L_3+L_2)+2(L_0+L_1)\\
  &=\frac{1}{3}(y(-1)+4y(0)+y(1))
\end{align*}
It is the same as Simpson's rule.
\subsection*{b}
\begin{align*}
  E^s(y) &= -\int_{-1}^{1} y(x) \,dx - \int_{-1}^{1} p_3(x) \,dx\\
  &=-\int_{-1}^{1} \frac{y^{(4)}(\xi )}{4!}(x+1)x^2(x-1),dx \\
  &=-\frac{y^{(4)}(\xi )}{90}
\end{align*}
\subsection*{c}

$$\int_{a}^{b} y(x) \,dx \approx \frac{h}{3}[y(a)+y(b)+4\sum_{i = 1}^{n}y(x_{2i-1})+2\sum_{i = 1}^{n-1}y(x_{2i})]$$
$$E^s(y)= -y^{(4)}(\xi)\frac{(b-a)}{2880}(2h)^4=-\frac{b-a}{180}h^4y^{(4)}(\eta )$$


\section*{II}
$|f''(x) = e^{-x^2}(4x^2-2)|\leq 2$,$|f^{(4)}(x)=4e^{-x^2}(4x^4-12x^2+3)|\leq 12$.
\subsection*{a}

$$\int_{-1}^{1} f(x) \,dx = \frac{1}{2}\sum_{k = 0}^{n-1} [f(x_k)+f(x_{k+1})]-\frac{b-a}{12}h^2f''(\eta)$$
We already know that $f''(x) = e^{-x^2}(4x^2-2)$, so $|f''(\eta)|\leq 2$.
\begin{align*}
  \frac{h^2}{6}&<0.5*10^{-6}\\
  n^2&>\frac{1}{3}*10^6\\
  n&>577
\end{align*}

\subsection*{b}
From I, we have Simpson rule.
We already know that $|f^{(4)}(x)=4e^{-x^2}(4x^4-12x^2+3)|\leq 12$
\begin{align*}
  \frac{h^4}{180}12&<0.5*10^{-6}\\
  n^4&>\frac{2}{15}*10^6\\
  n&>19
\end{align*}

\section*{III}
\subsection*{a}
$\pi_2(t)$ is orthogonal to $\pi_1(t) \in \mathbb{P}^{1}$. 
\begin{align*}
  \int_{0}^{\infty } (c_0+c_1t)(t^2+at+b)e^{-t} \,dt &=  \int_{0}^{\infty } [(c_1t^3+)+(c_0+c_1a)t^2+(ac_0+c_1b)t+c_0b]e^{-t} \,dt
  & = 6c_1+2c_1a+2c_0+ac_0+c_1b+c_0b
\end{align*}
\begin{align*}
  \left\{
  \begin{array}{l}
  2 + a + b = 0\\
  6 + 2a + b = 0
  \end{array}
  \right.
\end{align*}
So, $a = -4, b = 2$, which means $\pi_2(t) = t^2-4t+2$.
\subsection*{b}
$\pi_2(t) = t^2-4t+2 = 0$,we get $t_1 = 2 - \sqrt{2},t_2 = 2 + \sqrt{2}$.
\begin{align*}
  \left\{
  \begin{array}{l}
  \int_{0}^{\infty} e^{-x} \,dx = 1 = A_1 + A_2\\
  \int_{0}^{\infty} te^{-x} \,dx = 1 = A_1(2-\sqrt{2})+A_2(2+\sqrt{2}) 
  \end{array}
  \right.
\end{align*}
So, $A_1 = \frac{2+\sqrt{2}}{4},A_2 = \frac{2-\sqrt{2}}{4}$ and $E_2(f) = \frac{f^{(4)}(\xi)}{4!} \int_{0}^{\infty} \pi_2^2(x)e^{-x} \,dx = \frac{f^{(4)}(\xi)}{6}$.

\subsection*{c}
We have $f(x) = \frac{1}{1+x}$,$f(x_1)=\frac{3+\sqrt{2}}{7}$,$f(x_2)=\frac{3-\sqrt{2}}{7}$.

  $$\int_{0}^{\infty} \frac{1}{1+x}e^{-x} \,dx = \frac{3+\sqrt{2}}{7}A_1+\frac{3-\sqrt{2}}{7}A_2=\frac{4}{7}$$
$E_2(f)= I - \frac{4}{7}$, so we can get $24(1+\tau )^{-5}=6(I-\frac{4}{7})$,$\tau = 1.0926$

\section*{IV}
\subsection*{a}
The basic Lagrange polynomial \( \ell_m(t) \) is defined as:

\[
\ell_m(t) = \prod_{\substack{j=1 \\ j \neq m}}^{n} \frac{t - x_j}{x_m - x_j}.
\]
According to the problem, \( h_m(t) \) and \( q_m(t) \) are of the form:

\[
h_m(t) = (a_m + b_m t) \ell_m^2(t), \quad q_m(t) = (c_m + d_m t) \ell_m^2(t).
\]
To satisfy the Hermite interpolation conditions \( p(x_m) = f_m \) and \( p'(x_m) = f'_m \), we require:

\[
h_m(x_m) = 1, \quad h'_m(x_m) = 0.
\]

Substituting \( h_m(t) = (a_m + b_m t) \ell_m^2(t) \), we obtain:

\[
h_m(x_m) = a_m + b_mx_m = 1,
\]
\[
h'_m(x_m) = (b_m + 2a_m \ell'_m(x_m)) \ell_m^2(x_m) = b_m + 2a_m \ell'_m(x_m) = 0.
\]

Since \( \ell'_m(x_m) = \frac{d}{dt} \ell_m(t) \bigg|_{t=x_m} \) and \( \ell_m(x_m) = 1 \), we have:

\[
b_m = -2a_m \ell'_m(x_m) = -2 \ell'_m(x_m).
\]

Thus, the constants for \( h_m(t) \) are:

\[
a_m = 1 + 2x_m\ell'_m(x_m), \quad b_m = -2 \ell'_m(x_m).
\]
Similarly, to satisfy the Hermite interpolation conditions, we require:

\[
q_m(x_m) = 0, \quad q'_m(x_m) = 1.
\]

Substituting \( q_m(t) = (c_m + d_m t) \ell_m^2(t) \), we obtain:

\[
q_m(x_m) = c_m +d_mx_m = 0,
\]
\[
q'_m(x_m) = (d_m + 2c_m \ell'_m(x_m)) \ell_m^2(x_m) = d_m + 2c_m \ell'_m(x_m) = d_m = 1.
\]

Thus, the constants for \( q_m(t) \) are:

\[
c_m = -x_m, \quad d_m = 1.
\]\cite*{deepseek}

\subsection*{b}
\begin{align*}
  \left\{
  \begin{array}{l}
    \omega_k = \int_{a}^{b} (a_k+b_kx)l_k^2(x)\rho(x)  \,dx \\
    \mu_k = \int_{a}^{b} (c_k+d_kx)l_k^2(x)\rho(x)  \,dx
  \end{array}
  \right.
\end{align*}

\subsection*{c}
$$\mu_k = \int_{a}^{b} (c_k+d_kx)l_k^2(x)\rho(x)  \,dx = \int_{a}^{b} (x - x_k)l_k^2(x)\rho(x)  \,dx =0  $$
To ensure that \( \mu_k = 0 \) for each \( k = 1, 2, \ldots, n \), the following conditions must be imposed on the node polynomial or the nodes \( x_k \):

1. \textbf{Symmetry Condition}: If the nodes \( x_k \) are symmetrically distributed around a point (e.g., \( x_k = -x_{n-k+1} \)), then the integral of \( q_k(t) \) will be zero. This is because the term \( q_k(t) = (c_k + d_k t) \ell_k^2(t) \) will integrate to zero due to symmetry.

2. \textbf{Orthogonality Condition}: The node polynomial \( \pi(t) \) must be orthogonal to the polynomial \( t \ell_k^2(t) \) with respect to the weight function \( \rho(t) \). This ensures that the integral \( \int_{a}^{b} t \ell_k^2(t) \rho(t) \, dt = 0 \), which implies \( \mu_k = 0 \).

3. \textbf{Specific Node Distribution}: If the nodes \( x_k \) are chosen as the roots of a specific orthogonal polynomial (e.g., Legendre or Laguerre polynomials), then \( \mu_k = 0 \) is automatically satisfied due to the properties of these polynomials.


% ===============================================
\section*{ \center{\normalsize {Acknowledgement}} }
\printbibliography


\end{document}